\documentclass{article}
\usepackage[T1]{fontenc}
\usepackage{lmodern}
\usepackage{graphicx}
\usepackage{amsmath,amssymb}
\usepackage[a4paper,margin=2cm]{geometry}
\usepackage[absolute,overlay]{textpos}
\setlength{\TPHorizModule}{1mm}
\setlength{\TPVertModule}{1mm}
\usepackage[indonesian]{babel}
\usepackage{hyperref}
\setlength{\emergencystretch}{2em}
% unit: Halaman 4

\begin{document}

% === Halaman 4 (llm) ===

\section*{Grup}

\begin{itemize}
    \item Grup (\textit{group}) adalah sistem aljabar yang terdiri dari:
    \begin{itemize}
        \item sebuah himpunan $G$
        \item sebuah operasi biner $*$
    \end{itemize}
    sedemikian sehingga untuk semua elemen $a, b$, dan $c$ di dalam $G$ berlaku aksioma berikut:
    \begin{enumerate}
        \item \textit{Closure}: $a * b$ harus berada di dalam $G$
        \item Asosiatif: $a * (b * c) = (a * b) * c$
        \item Elemen netral: terdapat $e \in G$ sedemikian sehingga $a * e = e * a = a$
        \item Elemen invers: terdapat $a' \in G$ sedemikian sehingga $a * a' = a' * a = e$
    \end{enumerate}
    \item Notasi: $\langle G, *\rangle$
\end{itemize}

\end{document}
